% !TEX program = xelatex
\documentclass[a4paper,12pt]{article}

% ==================== 中文支持 ====================
\usepackage[UTF8]{ctex}

% ==================== 页面设置 ====================
\usepackage{geometry}
\geometry{left=2.5cm,right=2.5cm,top=2.5cm,bottom=2.5cm}

% ==================== 数学公式 ====================
\usepackage{amsmath}
\usepackage{amsfonts}
\usepackage{amssymb}
\usepackage{bm}

% ==================== 图表 ====================
\usepackage{graphicx}
\usepackage{float}
\usepackage{caption}
\captionsetup{labelsep=quad}

% ==================== 表格 ====================
\usepackage{array}
\usepackage{multirow}
\usepackage{booktabs}
\usepackage{longtable}

% ==================== 代码 ====================
\usepackage{listings}
\usepackage{xcolor}
\lstset{
	basicstyle=\small\ttfamily,
	breaklines=true,
	frame=single,
	numbers=left,
	numberstyle=\tiny,
	keywordstyle=\color{blue},
	commentstyle=\color{green!60!black},
	stringstyle=\color{red},
	showstringspaces=false
}

% ==================== 其他 ====================
\usepackage{hyperref}
\usepackage{indentfirst}
\setlength{\parindent}{2em}


% ==================== 标题格式（x.x.x 格式） ====================
\usepackage{titlesec}

% 一级标题：1、2、3、...
\titleformat{\section}{\centering\Large\bfseries}{\zhnum{section}、}{0.5em}{}

% 二级标题：1.1、1.2、...
\titleformat{\subsection}{\large\bfseries}{\thesubsection}{0.5em}{}

% 三级标题：1.1.1、1.1.2、...
\titleformat{\subsubsection}{\normalsize\bfseries}{\thesubsubsection}{0.5em}{}

% 四级标题：（1）（2）（3）...
\titleformat{\paragraph}{\normalsize\bfseries}{（\arabic{paragraph}）}{0.5em}{}



% ==================== 正文开始 ====================
\begin{document}
	
	% ==================== 摘要 ====================
	\setcounter{page}{1}
	\begin{center}
		{\LARGE\bfseries 基于[模型名称]的[问题名称]研究} \\[0.8cm]
		{\large 摘要} \\[0.3cm]
	\end{center}
	
	\noindent [研究背景简述]。本文针对[问题情境]建立了数学模型进行分析。\\[0.3cm]
	
	\noindent \textbf{针对问题一}，[方法简述]，得到[结论]。\\[0.3cm]
	
	\noindent \textbf{针对问题二}，[方法简述]，得到[结论]。\\[0.3cm]
	
	\noindent \textbf{针对问题三}，[方法简述]，得到[结论]。\\[0.3cm]
	
	\noindent \textbf{针对问题四}，[方法简述]，得到[结论]。\\[0.3cm]
	
	\noindent 最后，对模型进行了评价与推广。\\[0.3cm]
	
	\noindent \textbf{关键词：} [关键词1]；[关键词2]；[关键词3]；[关键词4]
	
	\newpage
	
	% ==================== 一、问题重述 ====================
	\section{问题重述}
	
	\subsection{问题背景}
	[内容]
	
	\subsection{问题要求}
	\textbf{问题一：} [描述]\\
	\textbf{问题二：} [描述]\\
	\textbf{问题三：} [描述]\\
	\textbf{问题四：} [描述]
	
	\newpage
	
	% ==================== 二、问题分析 ====================
	\section{问题分析}
	
	\subsection{问题一的分析}
	[内容]
	
	\subsection{问题二的分析}
	[内容]
	
	\subsection{问题三的分析}
	[内容]
	
	\subsection{问题四的分析}
	[内容]
	
	\newpage
	
	% ==================== 三、模型假设 ====================
	\section{模型假设}
	
	\begin{enumerate}
		\item 假设[内容]。
		\item 假设[内容]。
		\item 假设[内容]。
	\end{enumerate}
	
	\newpage
	
	% ==================== 四、符号说明 ====================
	\section{符号说明}
	
	\begin{table}[H]
		\centering
		\begin{tabular}{c|c|c}
			\hline
			\textbf{符号} & \textbf{含义} & \textbf{单位} \\
			\hline
			$x_i$ & [含义] & [单位] \\
			$P_i$ & [含义] & [单位] \\
			$S_{ij}$ & [含义] & [单位] \\
			\hline
		\end{tabular}
	\end{table}
	
	\newpage
	
	% ==================== 五、数据侧写 ====================
	\section{数据侧写}
	
	\subsection{[维度一]}
	[内容]，如图1所示。
	
	\begin{figure}[H]
		\centering
		\includegraphics[width=0.7\textwidth]{example-image}
		\caption{[图标题]}
		\label{fig:1}
	\end{figure}
	
	由图1可知，[结论]。
	
	\newpage
	
	% ==================== 六、模型的建立与求解 ====================
	\section{模型的建立与求解}
	
	% ----------------------------------------------
	% 问题一
	% ----------------------------------------------
	\subsection{问题一模型的建立与求解}
	
	\subsubsection{数据预处理}
	[内容]
	
	\subsubsection{模型建立}
	
	% 三级标题：1、2、3、...
	\paragraph{（1）决策变量的确定}
	[内容]
	
	\paragraph{（2）目标函数的确定}
	[公式]
	
	\paragraph{（3）约束条件的设定}
	[公式]
	
	\subsubsection{模型求解}
	
	\paragraph{（1）求解方法}
	[内容]
	
	\paragraph{（2）求解结果}
	[表格/图表]
	
	\subsubsection{结果分析}
	[内容]
	
	% ----------------------------------------------
	% 问题二
	% ----------------------------------------------
	\subsection{问题二模型的建立与求解}
	
	\subsubsection{模型建立}
	
	\paragraph{（1）不确定性因素分析}
	[内容]
	
	\paragraph{（2）目标函数的修正}
	[公式]
	
	\paragraph{（3）约束条件的补充}
	[公式]
	
	\subsubsection{模型求解}
	
	\paragraph{（1）求解算法}
	[内容]
	
	\paragraph{（2）求解结果}
	[表格/图表]
	
	\subsubsection{结果分析}
	[内容]
	
	% ----------------------------------------------
	% 问题三
	% ----------------------------------------------
	\subsection{问题三模型的建立与求解}
	
	\subsubsection{模型建立}
	[内容]
	
	\subsubsection{模型求解}
	[内容]
	
	\subsubsection{结果分析}
	[内容]
	
	% ----------------------------------------------
	% 问题四
	% ----------------------------------------------
	\subsection{问题四模型的建立与求解}
	
	\subsubsection{模型建立}
	[内容]
	
	\subsubsection{模型求解}
	[内容]
	
	\subsubsection{结果分析}
	[内容]
	
	\newpage
	
	% ==================== 七、AI工具使用声明 ====================
	\section*{AI工具使用声明}
	
	% 二选一：
	
	% 选项一：未使用
	本参赛队在竞赛过程中未使用任何AI工具。
	
	% 选项二：使用了（删除选项一，启用此项）
	% 本参赛队在竞赛过程中使用了AI工具，主要用于[简要用途]，详细使用情况见支撑材料。
	
	\vspace{1cm}
	\noindent 参赛队成员签名：\hspace{3cm} 日期：\hspace{3cm}
	
	\newpage
	
	% ==================== 八、参考文献 ====================
	\section{参考文献}
	
	\begin{enumerate}
		\item 作者1，作者2. 论文题目[J]. 期刊名称，年份，卷号(期号)：页码-页码.
		\item 作者. 学位论文题目[D]. 学校名称，年份.
		\item 作者. 书籍题目[M]. 出版社，年份.
	\end{enumerate}
	
	\newpage
	
	% ==================== 九、模型的评价与推广 ====================
	\section{模型的评价与推广}
	
	\subsection{模型的优点}
	\begin{enumerate}
		\item [优点1]
		\item [优点2]
	\end{enumerate}
	
	\subsection{模型的缺点}
	\begin{enumerate}
		\item [缺点1]
		\item [缺点2]
	\end{enumerate}
	
	\subsection{模型的推广}
	[内容]
	
	\newpage
	
	% ==================== 十、附录 ====================
	\section{附录}
	
	\subsection{支撑材料文件列表}
	\begin{table}[H]
		\centering
		\begin{tabular}{c|c}
			\hline
			\textbf{文件分类} & \textbf{文件内容} \\
			\hline
			问题一 & [文件列表] \\
			问题二 & [文件列表] \\
			问题三 & [文件列表] \\
			问题四 & [文件列表] \\
			\hline
		\end{tabular}
	\end{table}
	
	\subsection{核心代码}
	\begin{lstlisting}[language=Matlab, basicstyle=\small\ttfamily, caption=核心算法代码]
		% [代码内容]
	\end{lstlisting}
	
\end{document}